\subsection{赋范代数中的可乘序列}

设 A 是非离散赋值域 K 上的一个赋范代数（第 IX 章，§ 3，第 7 号，定义 9）；我们用 \(||x||\) 表示元素 \(x \in A\) 的范数，并设此范数满足不等式 \(||xy|| \leqslant ||x|| \cdot ||y||\) ；又设 A 有单位元 \(e\) 。设 \((x_n)_{n \in \mathbb{N}}\) 是 A 的点的一个无穷序列。N 的每个有限子集 J （以 N 的序线性排序）定义 A 的点的一个序列 \((x_n)_{n \in J}\) ，而我们定义积

\[
p _ {J} = \prod_ {n \in J} x _ {n}
\]

为此序列的积；此积称为序列 \((\pmb{x}_n)_{n\in \mathbb{N}}\) 的对应于有限子集 \(J\) （ \(\mathbb{N}\) 的）的有限部分积（回忆若 \(J = \emptyset\) ，我们令 \(\prod_{n\in \emptyset}x_n = e\) ）。

定义 1. 序列 \((\pmb{x}_n)_{n\in \mathbb{N}}\) 称为在赋范代数 \(\mathbf{A}\) 中可乘的，如果映射 \(\mathbf{J}\to \pmb{p}_{\mathbf{J}}\) 关于有限子集所成的集 \(\mathfrak{F}(\mathbf{N})\) （ \(\mathbf{N}\) 的）的截口滤有极限，此集以关系 \(c\) 排序；此极限称为序列 \((\pmb{x}_n)_{n\in \mathbb{N}}\) 的积，记作 \(\prod_{n\in \mathbb{N}}\pmb{x}_n\) （或简作 \(\prod_{n}\pmb{x}_n\) ）；诸 \(\pmb{x}_n\) 称为此积的因子。

定义 1 等价于下述：序列 \((\pmb{x}_n)\) 可乘且其积为 \(\pmb{p}\) ，如果对每个 \(\varepsilon > 0\) 存在有限子集 \(J_0\) （ \(\mathbf{N}\) 的），使得对每个有限子集 \(J \supset J_0\) （ \(\mathbf{N}\) 的），有 \(||\pmb{p}_{\mathrm{J}} - \pmb{p}|| \leqslant \varepsilon\) 。

注。1) 当 A 是交换代数时，定义 1 与第 III 章，§ 5，第 1 号，注 3 中所给的相同；但当 A 不交换时，指标集 N 的序结构本质上涉及定义 1 。若 \(\sigma\) 是 N 的任意置换，一般不能断言序列 \((\pmb{x}_{\sigma(n)})\) 在序列 \((x_{n})\) 可乘时可乘；且若两个序列都可乘，它们的积一般不同。

\begin{enumerate}
\def\labelenumi{\arabic{enumi})}
\setcounter{enumi}{1}
\tightlist
\item
  定义 1 可以直接推广到族 \((\pmb{x}_n)_{n\in I}\) 的情形，其指标集 I 是 \(\mathbf{Z}\) 的子集（以 \(\mathbf{Z}\) 的序所诱导的序线性排序）。我们将下述结果向此情形的推广留给读者（参看习题 1 与习题 2）。
\end{enumerate}

